\Askhsh{26183}{4}
\begin{ylh}{1}
    \item 1.5 Ιδιότητες των ορίων
    \item 1.8 Συνέχεια συνάρτησης
    \item 2.1 Η έννοια της παραγώγου
    \item 2.3 Κανόνες παραγώγισης
    \item 2.7 Τοπικά ακρότατα συνάρτησης
    \item 3.5 Η συνάρτηση [ορισμένο ολοκλήρωμα της f από α έως χ]
    \item 3.7 Εμβαδόν επιπέδου χωρίου.
\end{ylh}
Θεωρούμε τη συνάρτηση $f(x)$ με τύπο:
\[ f(x) = \begin{cases} \varepsilon \tan\left(\frac{\pi x}{4}\right), & 0 \le x \le 1 \\ 1 - \frac{4\ln x}{x}, & x > 1 \end{cases} \]
\begin{alist}
\item Να αποδείξετε ότι η $f$ είναι συνεχής στο $1$, αλλά όχι παραγωγίσιμη στο $1$. \monades{8}
\item Να αποδείξετε ότι η $f$ έχει ακριβώς δύο κρίσιμα σημεία στο διάστημα $[0, +\infty)$. \monades{7}
\item Να αποδείξετε ότι το εμβαδόν του χωρίου που ορίζεται από την γραφική παράσταση της $f$, τον άξονα $x'x$, τον άξονα $y'y$ και την ευθεία με εξίσωση $x = 1$, είναι $E = \frac{\ln 4}{\pi}$ τετραγωνικές μονάδες. \monades{10}
\end{alist}